\begin{proofbox}
	\textbf{\textcolor{CProof}{思路提示}}
	利用圆的切线垂直于过切点的半径这一几何性质，将垂直条件转化为坐标关系，直接导出切线方程。

	\medskip
	\textbf{\textcolor{CProof}{正式推导}}
	\begin{description}[style=nextline, leftmargin=2.8em, labelsep=0.5em, itemsep=0.55em]
		\item[\textbf{步骤一（垂直条件）：}]
			设圆心为 $O(0,0)$，切点为 $P(x_0,y_0)$，则半径 $OP$ 垂直于切线，因此向量 $\vec{OP}=(x_0,y_0)$ 是切线的一个法向量。
			\[
				\vec{n} = (x_0, y_0)
			\]
			。
			\par\textbf{理由：} 圆上任一点的切线与该点半径互相垂直，这是圆的几何性质。
		\item[\textbf{步骤二（点法式方程）：}]
			切线过点 $P(x_0,y_0)$，法向量为 $\vec{n}=(x_0,y_0)$，故点法式方程为
			\[
				x_0(x-x_0) + y_0(y-y_0) = 0
			\]
			。
			\par\textbf{理由：} 点法式直线方程 $a(x-x_0)+b(y-y_0)=0$ 中 $(a,b)$ 为法向量坐标。
		\item[\textbf{步骤三（化简并代换）：}]
			展开并移项，同时利用点 $P$ 在圆上 $x_0^2+y_0^2=r^2$，得
			\[
				\begin{aligned}
					x_0(x-x_0)+y_0(y-y_0) &= 0 \\ x_0 x - x_0^2 + y_0 y - y_0^2 &= 0 \\ x_0 x + y_0 y &= x_0^2 + y_0^2 \\ &= r^2
			\end{aligned}
			\]
			。
			\par\textbf{理由：} 前三行为代数变形，最后一行为代入 $x_0^2+y_0^2=r^2$ 得到切线方程标准形式。
	\end{description}

	\medskip
	\textbf{\textcolor{CProof}{结论回扣}}
	因此，过圆 $x^2+y^2=r^2$ 上一点 $P(x_0,y_0)$ 的切线方程为 $x_0 x + y_0 y = r^2$。该过程体现了“垂直$\to$法向量$\to$点法式$\to$化简”的完整思路。
\end{proofbox}
